\chapter{Tree Decompositions and Treewidth}

In the preceding chapters, we saw how trees yield fast queries and clean dynamic programs
because every edge is a cut and ancestors order subproblems neatly. But most networks
contain cycles. Tree decompositions carry the tree's hierarchical structure over to
such graphs. They measure, with an integer called
\emph{treewidth}, how far a graph deviates from a tree structure -- and turn that
measurement into an algorithm. A graph of treewidth $k$ inherits tree-like dynamic
programming, paying an overhead that depends on $k$ rather than $n$.

This is the structural foundation of \emph{fixed-parameter tractability}, the subject
of the next chapter. Here we build the machinery: what a tree decomposition is, why
the two technical conditions on it are exactly the right ones, how to run dynamic
programming over a decomposition (we will trace independent set and dominating set
in full, by hand, on a small graph), and how to compute or approximate treewidth
itself. By the end you should recognize a low-treewidth graph as one that splits into
small, nearly independent pieces that a dynamic program can solve efficiently.

\begin{historybox}
Tree decompositions were introduced in 1984 by Neil Robertson and Paul Seymour in Paper~II
of their landmark twenty-paper \emph{Graph Minors} project, which ultimately proved Wagner's
Conjecture: every minor-closed family of finite graphs is characterized by a finite set of
forbidden minors (where a graph $H$ is a minor of $G$ if $H$ can be formed by deleting
vertices, deleting edges, and contracting edges). Independently, related ideas surfaced in
Umberto Bertel{\'e} and Francesco Brioschi's work on non-serial dynamic programming (1972),
Rudolf Halin's study of $S$-functions (1976), and the partial $k$-trees of Stefan Arnborg and
Andrzej Proskurowski (1989). Robertson and Seymour needed a structural invariant to show that
any graph avoiding a fixed planar minor decomposes into tree-like pieces -- the celebrated
Excluded Grid Theorem.

Early attempts to extend tree algorithms to graphs centered on deleting feedback vertex
sets (sets of vertices whose removal leaves an acyclic graph). But that strategy is brittle:
an $n\times n$ grid requires deleting $\Th(n^2)$ vertices to kill all cycles, destroying all
local structure even though the grid is locally simple. Recursive partitioning also struggled,
because unconstrained cuts caused dynamic programming states to blow up.

Robertson and Seymour's insight was that the decomposition itself must form an abstract
tree of overlapping vertex subsets (``bags''). The critical condition is the \emph{running
intersection property} (every vertex induces a connected subtree of bags), guaranteeing that
bags act as separators that isolate past choices from future ones. In 1990, Bruno Courcelle
proved that every graph property expressible in Monadic Second-Order Logic ($\mathrm{MSO}_2$,
logic quantifying over vertices, edges, and sets of them) can be decided in linear time
$f(k)\cdot n$ on graphs of bounded treewidth $k$. In 1996, Hans Bodlaender showed that
deciding whether $\tw(G)\le k$ and finding an optimal decomposition is also solvable in
linear time $f(k)\cdot n$ for fixed $k$. Whether treewidth can be approximated within a
constant factor in polynomial time remains a central open problem in graph algorithms.
\end{historybox}

\chapterline

\section{From Trees to Tree-Like Graphs}
\label{ch24:intro}

\begin{intuitionbox}
A tree is easy to work with for two reasons. First, a tree has tiny \emph{separators}:
delete one vertex and the tree falls into independent pieces, so a subproblem in
one piece interacts with the rest only through that single cut vertex. Second,
those separators are organized hierarchically -- the recursion tree of ``cut
here, recurse on the pieces'' is itself a tree, so a dynamic program can sweep it
once. A tree decomposition provides both properties for a general graph. It
groups the vertices into overlapping sets called \emph{bags}, arranges the bags
in a tree, and guarantees that each bag acts as a small separator between the
parts of the graph on either side of it. The largest bag, minus one, is the
treewidth, which governs the cost of tree-like dynamic programming.
\end{intuitionbox}

Consider the cycle $C_4$ on vertices $\{1,2,3,4\}$ with edges $12,23,34,41$. It
is not a tree -- it has a cycle -- but it is barely not a tree. Cut it at any two
opposite vertices and it falls apart. We can capture ``barely not a tree'' by
covering the graph with bags of size three, arranged in a path:

\begin{center}
\begin{tikzpicture}[>=Stealth, every node/.style={font=\small}]
  % the graph C4
  \begin{scope}[xshift=0cm]
    \node[circle,draw=nodeblue,fill=blue!5,thick,minimum size=7mm] (v1) at (0,1.4) {$1$};
    \node[circle,draw=nodeblue,fill=blue!5,thick,minimum size=7mm] (v2) at (1.4,1.4) {$2$};
    \node[circle,draw=nodeblue,fill=blue!5,thick,minimum size=7mm] (v3) at (1.4,0) {$3$};
    \node[circle,draw=nodeblue,fill=blue!5,thick,minimum size=7mm] (v4) at (0,0) {$4$};
    \draw[thick] (v1)--(v2)--(v3)--(v4)--(v1);
    \node at (0.7,-0.9) {the graph $C_4$};
  \end{scope}
  % the decomposition
  \begin{scope}[xshift=6cm]
    \node[rounded corners,draw=structurepurple,fill=purple!6,thick,minimum height=9mm] (b1) at (0,0.8) {$\{1,2,3\}$};
    \node[rounded corners,draw=structurepurple,fill=purple!6,thick,minimum height=9mm] (b2) at (0,-0.8) {$\{1,3,4\}$};
    \draw[thick] (b1)--(b2);
    \node at (0,-1.9) {a tree decomposition, width $2$};
  \end{scope}
\end{tikzpicture}
\end{center}

Two bags, $\{1,2,3\}$ and $\{1,3,4\}$, joined by an edge. Every vertex appears in
some bag. Every edge of $C_4$ -- including the pairs $12,23$ in the first bag and
$34,41$ in the second -- sits inside a single bag. And the shared pair $\{1,3\}$
is exactly the separator that breaks the cycle. The widest bag has three vertices,
so the \emph{width} is $3-1=2$. We will prove shortly that $C_4$ has treewidth
exactly $2$, and more generally that trees are precisely the graphs of treewidth
$1$.

\section{The Definition and Why Each Clause Matters}
\label{ch24:defs}

\begin{definition}[Tree decomposition]
\label{ch24:def-td}
A \emph{tree decomposition} of a graph $\Gph=(V,E)$ is a pair $(\mathcal{T},
\{B_x\}_{x\in \mathcal{T}})$ where $\mathcal{T}$ is a tree and each node $x$ of
$\mathcal{T}$ carries a \emph{bag} $B_x\subseteq V$, subject to three conditions:
\begin{itemize}
  \item[\textbf{(V)}] \emph{Vertex cover:} $\bigcup_{x\in\mathcal{T}} B_x = V$.
        Every vertex lives in at least one bag.
  \item[\textbf{(E)}] \emph{Edge cover:} for every edge $\{u,v\}\in E$ there is a
        node $x$ with $\{u,v\}\subseteq B_x$. Every edge is realized inside some
        bag.
  \item[\textbf{(C)}] \emph{Connectivity (coherence):} for every vertex $v\in V$,
        the set of nodes $\{x : v\in B_x\}$ induces a \emph{connected subtree} of
        $\mathcal{T}$.
\end{itemize}
The \emph{width} of the decomposition is $\max_{x}\abs{B_x}-1$. The
\emph{treewidth} $\tw(\Gph)$ is the minimum width over all tree decompositions.
\end{definition}

The ``$-1$'' is a convention chosen so that trees come out with treewidth $1$
rather than $2$; see Proposition~\ref{ch24:prop-trees}. To keep the names
straight: nodes (of $\mathcal{T}$) carry bags; vertices (of $\Gph$) live inside
bags. We reserve ``node'' for $\mathcal{T}$ and ``vertex'' for $\Gph$.

\begin{focusbox}
Each of the three conditions has a purpose. \textbf{(V)} and \textbf{(E)} say
the decomposition does not lose information: you could reconstruct $\Gph$ from the
bags. The essential condition is \textbf{(C)}. It forces each vertex to occupy a
connected subtree of $\mathcal{T}$, and that connectedness is what makes bags act as
separators. Without \textbf{(C)}, a vertex could appear in two distant bags
without appearing along the path between them, creating dependencies across
subtrees and breaking dynamic programming. Condition \textbf{(C)} says each
vertex enters the decomposition once, stays present across a single connected
region of $\mathcal{T}$, and never reappears after leaving.
\end{focusbox}

\begin{figure}[H]
\centering
\begin{tikzpicture}[>=Stealth,
  vnode/.style={circle,draw=sourceblue,fill=sourceblue!12,thick,minimum size=7mm,inner sep=0pt,font=\small},
  bag/.style={rounded corners=3pt,draw=structurepurple,fill=structurepurple!8,thick,inner sep=5pt,font=\small}
]
  % Lane 1: The Underlying Graph G
  \begin{scope}[xshift=0cm]
    \node[font=\footnotesize\bfseries\scshape,text=sourceblue!80!black] at (1.5,2.4) {Underlying Graph $\Gph$};
    \node[vnode] (1) at (0,1.4) {1};
    \node[vnode] (2) at (1.5,1.4) {2};
    \node[vnode] (3) at (0,0) {3};
    \node[vnode] (4) at (1.5,0) {4};
    \node[vnode] (5) at (3.0,1.4) {5};
    \node[vnode] (6) at (3.0,0) {6};

    \draw[thick,sourceblue] (1)--(2) (1)--(3) (2)--(3) (2)--(4) (3)--(4) (2)--(5) (4)--(5) (4)--(6) (5)--(6);
    
    \node[font=\scriptsize,text width=36mm,align=center] at (1.5,-0.9) {6 vertices, 9 edges.\\Triangulated structure.};
  \end{scope}

  % Arrow between lanes
  \draw[->,very thick,structurepurple] (3.7,0.7) -- (5.1,0.7) node[midway,above=3pt,font=\scriptsize\bfseries] {decompose};

  % Lane 2: Tree of Bags
  \begin{scope}[xshift=6.6cm]
    \node[font=\footnotesize\bfseries\scshape,text=structurepurple!80!black] at (3.2,2.4) {Tree Decomposition $\mathcal{T}$ (Width $2$)};

    \node[bag] (b1) at (0,1.3)    {$B_1 = \{1,2,3\}$};
    \node[bag] (b2) at (3.3,1.3)  {$B_2 = \{2,3,4\}$};
    \node[bag] (b3) at (3.3,-0.3) {$B_3 = \{2,4,5\}$};
    \node[bag] (b4) at (6.6,-0.3) {$B_4 = \{4,5,6\}$};

    \draw[very thick,structurepurple] (b1) -- node[midway,above=2pt,font=\scriptsize] {$\{2,3\}$} (b2);
    \draw[very thick,structurepurple] (b2) -- node[midway,right=2pt,font=\scriptsize] {$\{2,4\}$} (b3);
    \draw[very thick,structurepurple] (b3) -- node[midway,above=2pt,font=\scriptsize] {$\{4,5\}$} (b4);

    % Callouts for coherence
    \node[font=\scriptsize,targetgreen!80!black,text width=74mm,align=left] at (3.3,-1.6) {
      \textbf{(V)} Every vertex in $\ge 1$ bag.\\
      \textbf{(E)} Every edge inside $\ge 1$ bag.\\
      \textbf{(C)} Subtree for vertex 2: $\{B_1, B_2, B_3\}$ (connected!).\\
      \textbf{(C)} Subtree for vertex 4: $\{B_2, B_3, B_4\}$ (connected!).
    };
  \end{scope}
\end{tikzpicture}
\caption{A tree decomposition of width $2$ on a $6$-vertex graph. Every edge of $\Gph$ is contained in at least one bag. The set of bags containing vertex $2$ forms a connected subpath $(B_1 - B_2 - B_3)$ in the tree, and similarly for vertex $4$ $(B_2 - B_3 - B_4)$. The intersection of adjacent bags $B_2\cap B_3 = \{2,4\}$ acts as a balanced separator cutting the left vertices $\{1,3\}$ from the right vertex $\{6\}$.}
\label{ch24:fig-bag-structure}
\end{figure}

\subsection{The separator property}

Every algorithm in this chapter rests on the following property.
Removing the vertices of a single bag disconnects the rest of the graph along the
tree.

\begin{lemma}[Bags are separators]
\label{ch24:lem-separator}
Let $(\mathcal{T},\{B_x\})$ be a tree decomposition of $\Gph$, and let
$\{x,y\}$ be an edge of $\mathcal{T}$. Deleting that edge splits $\mathcal{T}$
into two subtrees $\mathcal{T}_x\ni x$ and $\mathcal{T}_y\ni y$. Let
\[
  A = \Bigl(\bigcup_{z\in\mathcal{T}_x} B_z\Bigr)\setminus (B_x\cap B_y),
  \qquad
  C = \Bigl(\bigcup_{z\in\mathcal{T}_y} B_z\Bigr)\setminus (B_x\cap B_y).
\]
Then $S := B_x\cap B_y$ separates $A$ from $C$: no edge of $\Gph$ joins a vertex
of $A$ to a vertex of $C$.
\end{lemma}

\begin{proof}
Suppose for contradiction $\{a,c\}\in E$ with $a\in A$, $c\in C$. By \textbf{(E)}
some bag $B_z$ contains both $a$ and $c$. The node $z$ lies in $\mathcal{T}_x$ or
in $\mathcal{T}_y$; say $z\in\mathcal{T}_y$ (the other case is symmetric). Now
$a\in A$ means $a$ occurs in some bag of $\mathcal{T}_x$, and $a\in B_z$ means it
occurs in a bag of $\mathcal{T}_y$. By \textbf{(C)} the nodes containing $a$ form
a connected subtree, and any path in $\mathcal{T}$ between a node of
$\mathcal{T}_x$ and a node of $\mathcal{T}_y$ must cross the edge $\{x,y\}$. Hence
both $x$ and $y$ contain $a$, so $a\in B_x\cap B_y=S$. But $A$ was defined to
exclude $S$ -- contradiction.
\end{proof}

\begin{figure}[H]
\centering
\begin{tikzpicture}[>=Stealth, every node/.style={font=\small}]
  % left subtree T_x as a blob of A vertices
  \node[draw=nodeblue,thick,fill=blue!5,rounded corners,minimum width=22mm,minimum height=16mm] (A) at (0,0) {$A$};
  \node[font=\scriptsize] at (0,-1.3) {$\mathcal{T}_x$ side};
  % separator
  \node[draw=structurepurple,thick,fill=purple!10,rounded corners,minimum width=10mm,minimum height=16mm] (S) at (3,0) {$S$};
  \node[font=\scriptsize] at (3,-1.3) {$B_x\cap B_y$};
  % right subtree T_y as a blob of C vertices
  \node[draw=targetgreen,thick,fill=green!6,rounded corners,minimum width=22mm,minimum height=16mm] (C) at (6,0) {$C$};
  \node[font=\scriptsize] at (6,-1.3) {$\mathcal{T}_y$ side};
  % only edges allowed: A-S and S-C, never A-C
  \draw[thick] (A)--(S);
  \draw[thick] (S)--(C);
  \draw[trapred,thick,dashed] (A) to[bend left=35] node[midway,above,font=\scriptsize]{forbidden} (C);
\end{tikzpicture}
\caption{Lemma~\ref{ch24:lem-separator} in one picture. Deleting tree edge
$\{x,y\}$ partitions the graph vertices into $A$, the separator $S=B_x\cap B_y$,
and $C$. Edges may run $A$--$S$ or $S$--$C$, but never directly $A$--$C$: every
$A$--$C$ interaction is forced through $S$.}
\label{ch24:fig-separator}
\end{figure}

\begin{whybox}
Lemma~\ref{ch24:lem-separator} (Figure~\ref{ch24:fig-separator}) is the central
mechanism. A dynamic program processes the
tree one node at a time; when it crosses an edge of $\mathcal{T}$, everything it
needs to remember about the already-processed side is summarized by how the
solution behaves on the separator $S=B_x\cap B_y$, a set of at most $\tw+1$
vertices. The number of possible behaviors on $S$ is bounded by a function of
$\tw$ alone -- $2^{\abs S}$ for independent set, $3^{\abs S}$ for dominating set,
$k^{\abs S}$ for $k$-coloring. That is why small treewidth makes hard problems tractable:
the ``interface'' between subproblems is provably small.
\end{whybox}

\subsection{Trees, cycles, and cliques}

\begin{proposition}[Treewidth $1$ means forest]
\label{ch24:prop-trees}
$\tw(\Gph)\le 1$ if and only if $\Gph$ is a forest (acyclic). A graph with at
least one edge has $\tw\ge 1$.
\end{proposition}

\begin{proof}
($\Leftarrow$) Root each tree of the forest. Make one decomposition node per
\emph{edge} $\{u,\mathrm{parent}(u)\}$ with bag $\{u,\mathrm{parent}(u)\}$, and
hang these nodes in the shape of the tree (the node for $u$ is a child of the node
for $\mathrm{parent}(u)$); add a singleton bag for the root. Each bag has size
$\le 2$, so width $\le 1$. Conditions \textbf{(V)},\textbf{(E)} are immediate, and
\textbf{(C)} holds because $u$ appears only in its own edge-node and in those of
its children -- a connected subtree.

($\Rightarrow$) Take a decomposition of width $\le 1$, so every bag has at most
two vertices. We show $\Gph$ is a forest by induction on $\abs{V}$ plus the number
of tree nodes. If $T$ has one node, $\Gph$ has at most two vertices and is a
forest. Otherwise pick a leaf $x$ of $T$ with neighbor $y$. If
$B_x\subseteq B_y$, delete $x$; the rest is still a decomposition of $\Gph$.
Otherwise some vertex $v\in B_x$ is missing from $B_y$. Its subtree $T_v$ is
connected, contains the leaf $x$, and avoids $x$'s only neighbor $y$, so
$T_v=\{x\}$: the vertex $v$ lives in $B_x$ alone. Every edge at $v$ must sit in a
bag containing $v$, which means in $B_x$, which holds only one other vertex. So
$v$ has degree $\le 1$. Removing $v$ from $\Gph$ and from $B_x$ gives a width-$\le
1$ decomposition of $\Gph-v$, which is a forest by induction. Adding back a vertex
of degree $\le 1$ cannot close a cycle, so $\Gph$ is a forest. The last sentence
of the proposition holds because a single edge $\{u,v\}$ must, by
\textbf{(E)}, lie in a bag of size $\ge 2$.
\end{proof}

\begin{proposition}[Cliques are unavoidable]
\label{ch24:prop-clique}
If $K\subseteq V$ induces a clique in $\Gph$, then some bag contains all of $K$.
Consequently $\tw(\Gph)\ge \omega(\Gph)-1$, where $\omega$ is the clique number,
and $\tw(K_n)=n-1$.
\end{proposition}

\begin{proof}
This is the \emph{Helly property} for subtrees of a tree: a family of pairwise
intersecting subtrees has a common node. For each $v\in K$ let $T_v=\{x:v\in
B_x\}$, a subtree by \textbf{(C)}. For any two $u,v\in K$ the edge $\{u,v\}$ is in
some bag (clique, so the edge exists; \textbf{(E)}), hence $T_u\cap T_v\ne
\emptyset$ -- the subtrees pairwise intersect. The Helly property for subtrees of
a tree gives a common node $x\in\bigcap_{v\in K} T_v$, and then $K\subseteq B_x$.
So $\abs{B_x}\ge\abs K$, width $\ge\abs K-1$. For $K_n$ this forces a bag with all
$n$ vertices, width $n-1$; the one-bag decomposition achieves it.
\end{proof}

\section{A Worked Decomposition by Hand}
\label{ch24:worked-decomp}

Let us fix a running graph and find a good tree decomposition for it. Take
$\Gph$ on seven vertices $\{a,b,c,d,e,f,g\}$ with edges
\[
  ab,\; ac,\; bc,\; bd,\; cd,\; ce,\; de,\; ef,\; eg,\; fg.
\]

\begin{center}
\begin{tikzpicture}[>=Stealth, every node/.style={font=\small},
  V/.style={circle,draw=nodeblue,fill=blue!5,thick,minimum size=7mm}]
  \node[V] (a) at (0,1) {$a$};
  \node[V] (b) at (1.4,1.8) {$b$};
  \node[V] (c) at (1.4,0.2) {$c$};
  \node[V] (d) at (2.8,1.8) {$d$};
  \node[V] (e) at (3.4,0.2) {$e$};
  \node[V] (f) at (5.0,1.0) {$f$};
  \node[V] (g) at (5.0,-0.6) {$g$};
  \draw[thick] (a)--(b) (a)--(c) (b)--(c) (b)--(d) (c)--(d) (c)--(e) (d)--(e) (e)--(f) (e)--(g) (f)--(g);
\end{tikzpicture}
\end{center}

Here $\{a,b,c\}$, $\{b,c,d\}$, $\{c,d,e\}$ form a chain of
triangles sharing edges, then $\{e,f,g\}$ is a triangle hanging off $e$. This
``chain of overlapping triangles'' is exactly a graph of treewidth $2$. A natural
decomposition reads the chain off directly:

\begin{center}
\begin{tikzpicture}[>=Stealth, every node/.style={font=\small},
  Bag/.style={rounded corners,draw=structurepurple,fill=purple!6,thick,minimum height=8mm,inner sep=4pt}]
  \node[Bag] (x1) at (0,0) {$\{a,b,c\}$};
  \node[Bag] (x2) at (2.6,0) {$\{b,c,d\}$};
  \node[Bag] (x3) at (5.2,0) {$\{c,d,e\}$};
  \node[Bag] (x4) at (7.8,0) {$\{e,f,g\}$};
  \draw[thick] (x1)--(x2) node[midway,above,font=\scriptsize]{$b,c$};
  \draw[thick] (x2)--(x3) node[midway,above,font=\scriptsize]{$c,d$};
  \draw[thick] (x3)--(x4) node[midway,above,font=\scriptsize]{$e$};
\end{tikzpicture}
\end{center}

Check the conditions. \textbf{(V)}: all seven vertices appear. \textbf{(E)}: every
edge lives in a bag -- e.g.\ $ab,ac,bc\subseteq\{a,b,c\}$; $bd,cd\subseteq
\{b,c,d\}$; $ce,de\subseteq\{c,d,e\}$; $ef,eg,fg\subseteq\{e,f,g\}$. \textbf{(C)}:
trace each vertex. $c$ appears in the first three bags -- contiguous. $d$ in bags
$2,3$ -- contiguous. $e$ in bags $3,4$ -- contiguous. Good. Widest bag has $3$
vertices, so this decomposition has width $2$, and since $\Gph$ contains a
triangle (clique of size $3$), Proposition~\ref{ch24:prop-clique} gives
$\tw(\Gph)\ge 2$. Therefore $\tw(\Gph)=2$ exactly.

\begin{labbox}
Build a tree decomposition of the $3\times 3$ grid graph (vertices a $3\times 3$
lattice, edges between orthogonal neighbors). Try to get width $3$: use bags that
are ``two adjacent columns'' sliding across the grid, $\{\text{col }i,\text{col }
i{+}1\}$. Verify \textbf{(C)} holds for the middle column. Then convince yourself
you cannot do better than width $3$ -- the $3\times 3$ grid has treewidth exactly
$3$, and in general the $n\times n$ grid has treewidth $n$, the canonical example
of a graph that is sparse yet has unbounded treewidth.
\end{labbox}

\section{Nice Tree Decompositions}
\label{ch24:nice}

Dynamic programming wants a tidy tree where each step does exactly one small
thing. The raw decomposition above has bags changing by several vertices across a
single edge, which makes the transition rules messy. We standardize.

\begin{definition}[Nice tree decomposition]
\label{ch24:def-nice}
A tree decomposition is \emph{nice} if $\mathcal{T}$ is rooted, the root bag and
every leaf bag are empty, and every non-leaf node is one of:
\begin{itemize}
  \item \textbf{Introduce $v$}: one child $y$, and $B_x = B_y\cup\{v\}$ for a
        single $v\notin B_y$.
  \item \textbf{Forget $v$}: one child $y$, and $B_x = B_y\setminus\{v\}$ for a
        single $v\in B_y$.
  \item \textbf{Join}: two children $y,z$ with $B_x = B_y = B_z$.
\end{itemize}
\end{definition}

\begin{proposition}[Niceness is free]
\label{ch24:prop-nice}
Any width-$k$ tree decomposition with $\Oh(n)$ nodes can be converted to a nice
one of width $k$ with $\Oh(k\cdot n)$ nodes in $\Oh(k^2 n)$ time.
\end{proposition}

\begin{proof}[Proof sketch]
Root the tree arbitrarily. For an edge from parent $x$ to child $y$ where the bags
differ, insert a chain of introduce/forget nodes that adds the vertices of
$B_x\setminus B_y$ one at a time and removes those of $B_y\setminus B_x$ one at a
time; each step changes the bag by one vertex, at most $\abs{B_x}+\abs{B_y}\le
2(k+1)$ steps. For a node with $\ge 2$ children, replace it by a binary tree of
join nodes (all sharing the same bag), then attach the children. Finally pad the
top with forget nodes down to an empty root and cap each leaf with introduce nodes
from an empty bag. Every inserted bag is a subset of an existing one, so width is
unchanged; the total node count is $\Oh(kn)$.
\end{proof}

A nice tree decomposition can be drawn as a chain of one-vertex changes. Here is a nice
version of the bottom of the spine $\{a,b,c\}-\{b,c,d\}$ from our running graph:

\begin{center}
\begin{tikzpicture}[>=Stealth, every node/.style={font=\scriptsize},
  Bag/.style={rounded corners,draw=structurepurple,fill=purple!6,thick,minimum height=6.5mm,inner sep=3.5pt}]
  \node[Bag] (l)  at (0,0)    {$\emptyset$ (leaf)};
  \node[Bag] (i1) at (0,1.25) {$\{a\}$};
  \node[Bag] (i2) at (0,2.50) {$\{a,b\}$};
  \node[Bag] (i3) at (0,3.75) {$\{a,b,c\}$};
  \node[Bag] (f1) at (0,5.00) {$\{b,c\}$};
  \node[Bag] (i4) at (0,6.25) {$\{b,c,d\}$};
  \draw[->,thick] (l)--(i1)  node[midway,right=4pt]{intro $a$};
  \draw[->,thick] (i1)--(i2) node[midway,right=4pt]{intro $b$};
  \draw[->,thick] (i2)--(i3) node[midway,right=4pt]{intro $c$};
  \draw[->,thick] (i3)--(f1) node[midway,right=4pt]{forget $a$};
  \draw[->,thick] (f1)--(i4) node[midway,right=4pt]{intro $d$};
\end{tikzpicture}
\end{center}

The arrows point from child to parent (leaves at the bottom, root at top), the
direction the dynamic program flows.

\section{Dynamic Programming over a Tree Decomposition}
\label{ch24:dp}

The dynamic programming template is always the same: at each node $x$ we maintain a
table indexed by the possible ``states'' of the bag $B_x$, where the value stored
is the best (or the count, or feasibility) of a partial solution on the part of
the graph already absorbed below $x$, \emph{conditioned} on the bag being in that
state. Lemma~\ref{ch24:lem-separator} guarantees the bag's state is a complete
interface: nothing below $x$ talks to anything above $x$ except through $B_x$.

\subsection{Maximum weight independent set}

An \emph{independent set} is a set of pairwise non-adjacent vertices. Given vertex
weights $w(v)\ge 0$, we want a maximum-weight independent set (MWIS). NP-hard in
general; here is the bounded-treewidth algorithm.

For a node $x$ with bag $B_x$, a \emph{state} is a subset $S\subseteq B_x$ of bag
vertices we tentatively include, required to be independent inside the bag. Let
$V_x$ denote the union of all bags in the subtree rooted at $x$ (the vertices
``introduced so far''). Define
\[
  \mathrm{dp}[x][S] \;=\;
  \max\Bigl\{\; w(I) \;:\; I\subseteq V_x \text{ independent},\;
  I\cap B_x = S \;\Bigr\},
\]
the best weight of an independent set on $V_x$ whose intersection with the current
bag is exactly $S$ (and $-\infty$ if $S$ itself is not independent in $\Gph$). The
transition rules:

\begin{itemize}
  \item \textbf{Leaf} ($B_x=\emptyset$): the only state is $S=\emptyset$, with
        $\mathrm{dp}[x][\emptyset]=0$.
  \item \textbf{Introduce $v$} (child $y$): for each independent $S\subseteq B_x$,
        \[
          \mathrm{dp}[x][S]=
          \begin{cases}
            \mathrm{dp}[y][S] & v\notin S,\\[2pt]
            \mathrm{dp}[y][S\setminus\{v\}] + w(v) & v\in S,
          \end{cases}
        \]
        provided $S$ stays independent; if $v\in S$ but $S$ contains a neighbor of
        $v$, set $\mathrm{dp}[x][S]=-\infty$.
  \item \textbf{Forget $v$} (child $y$): for each $S\subseteq B_x$,
        \[
          \mathrm{dp}[x][S]=\max\bigl(\mathrm{dp}[y][S],\;
          \mathrm{dp}[y][S\cup\{v\}]\bigr).
        \]
        We project away $v$, keeping the better of ``$v$ out'' and ``$v$ in.''
  \item \textbf{Join} (children $y,z$, $B_x=B_y=B_z$): for each independent $S$,
        \[
          \mathrm{dp}[x][S]=\mathrm{dp}[y][S]+\mathrm{dp}[z][S]-w(S),
        \]
        where $w(S)=\sum_{v\in S} w(v)$. The two subtrees agree on $S$, so we add
        their values but subtract the weight of $S$ counted twice. For a
        non-independent $S$, set $\mathrm{dp}[x][S]=-\infty$, so every table
        entry is defined when the parent reads it.
\end{itemize}

The answer is $\mathrm{dp}[r][\emptyset]$ at the empty root $r$.

\begin{algorithm}[H]
\caption{\textsc{MWIS} on a nice tree decomposition}
\label{ch24:alg-mwis}
\KwIn{nice tree decomposition $(\mathcal{T},\{B_x\})$ rooted at $r$, weights $w$}
\KwOut{maximum weight of an independent set of $\Gph$}
\ForEach{node $x$ of $\mathcal{T}$ in post-order}{
  \uIf{$x$ is a leaf}{$\mathrm{dp}[x][\emptyset]\gets 0$\;}
  \uElseIf{$x$ introduces $v$ over child $y$}{
    \ForEach{$S\subseteq B_x$}{
      \uIf{$v\in S$ and $S$ is independent in $\Gph$}{
        $\mathrm{dp}[x][S]\gets \mathrm{dp}[y][S\setminus\{v\}]+w(v)$\;}
      \uElseIf{$v\notin S$}{$\mathrm{dp}[x][S]\gets \mathrm{dp}[y][S]$\;}
      \Else{$\mathrm{dp}[x][S]\gets -\infty$\;}
    }
  }
  \uElseIf{$x$ forgets $v$ over child $y$}{
    \ForEach{$S\subseteq B_x$}{
      $\mathrm{dp}[x][S]\gets\max\bigl(\mathrm{dp}[y][S],\,\mathrm{dp}[y][S\cup\{v\}]\bigr)$\;}
  }
  \Else(\tcp*[h]{join with children $y,z$}){
    \ForEach{$S\subseteq B_x$}{
      \uIf{$S$ is independent in $\Gph$}{
        $\mathrm{dp}[x][S]\gets \mathrm{dp}[y][S]+\mathrm{dp}[z][S]-w(S)$\;}
      \Else{$\mathrm{dp}[x][S]\gets -\infty$\;}
    }
  }
}
\Return $\mathrm{dp}[r][\emptyset]$\;
\end{algorithm}

\begin{whybox}
Why is the join rule subtract-the-overlap and not something fancier? Because the
separator $B_x$ splits $V_x$ into the $y$-side and the $z$-side, sharing only
$B_x$ itself (Lemma~\ref{ch24:lem-separator}). Fix the bag's chosen set $S$. Any
independent set on the whole subtree restricted to ``$\cap B_x=S$'' decomposes
uniquely into an independent set on the $y$-side and one on the $z$-side, each
intersecting the bag in $S$, and the separation guarantee ensures there are \emph{no edges} between the
private parts of the two sides to forbid -- so we may combine them freely. Their
weights add, but the bag vertices $S$ were paid for in both tables, hence
$-w(S)$. The independence \emph{within} $S$ was already enforced when $S$ was
introduced, and there are no other cross edges to check. That is the entire
content of the rule.
\end{whybox}

\begin{theorem}[Correctness of Algorithm~\ref{ch24:alg-mwis}]
\label{ch24:thm-mwis}
For every node $x$ and every $S\subseteq B_x$, Algorithm~\ref{ch24:alg-mwis}
computes $\mathrm{dp}[x][S]$ equal to the maximum weight of an independent set
$I\subseteq V_x$ with $I\cap B_x=S$ (or $-\infty$ if none exists). In particular
$\mathrm{dp}[r][\emptyset]$ is the MWIS weight of $\Gph$.
\end{theorem}

\begin{proof}
Induction on the post-order. \emph{Leaf:} $V_x=\emptyset$, only the empty set
qualifies, value $0$. \emph{Introduce $v$:} by \textbf{(C)}, $v$ is brand new
($v\notin V_y$ since $v\notin B_y$ and the subtree of $\mathcal{T}$ containing $v$
is rooted here), so every edge incident to $v$ within $V_x$ has its other endpoint
in $B_x$ -- this is the consequence of the separator property. If $v\in S$ we must
check $S$ is independent (done) and then $v$'s contribution is additive over the
child's best set with $v$ removed; if $v\notin S$, $v$ contributes nothing and any
$I$ counted by the child with $I\cap B_y=S$ still has $I\cap B_x=S$.
\emph{Forget $v$:} an independent set $I\subseteq V_x=V_y$ with $I\cap B_x=S$
either contains $v$ (then $I\cap B_y=S\cup\{v\}$) or not (then $I\cap B_y=S$);
taking the max over the two child entries is exactly right, and no new edges
appear because $V_x=V_y$. \emph{Join:} by Lemma~\ref{ch24:lem-separator}, uniqueness
of the split uses that the $y$-side private vertices and $z$-side private vertices
have no edge between them (separator), so the combined set is independent iff each
side is and they agree on $S$. Maximizing each side independently and adding
(correcting the double count) yields the optimum.
\end{proof}

\subsection{Tracing it by hand}

Take the path $P_3$: vertices $1-2-3$, edges $12,23$, all weights $w=1$. The MWIS
is $\{1,3\}$ with weight $2$. We run the DP on the nice decomposition with spine
that introduces $1$, introduces $2$, forgets $1$, introduces $3$, forgets $2$,
forgets $3$. We track the table as we climb.

\begin{center}
\renewcommand{\arraystretch}{1.25}
\begin{tabular}{l l l}
\toprule
node & bag & table (state $\mapsto$ value) \\
\midrule
leaf                & $\emptyset$ & $\emptyset\mapsto 0$ \\
intro $1$           & $\{1\}$     & $\emptyset\mapsto 0,\;\{1\}\mapsto 1$ \\
intro $2$           & $\{1,2\}$   & $\emptyset\mapsto 0,\;\{1\}\mapsto 1,\;\{2\}\mapsto 1,\;\{1,2\}\mapsto -\infty$ \\
forget $1$          & $\{2\}$     & $\emptyset\mapsto \max(0,1)=1,\;\{2\}\mapsto \max(1,-\infty)=1$ \\
intro $3$           & $\{2,3\}$   & $\emptyset\mapsto 1,\;\{3\}\mapsto 1{+}1{=}2,\;\{2\}\mapsto 1,\;\{2,3\}\mapsto -\infty$ \\
forget $2$          & $\{3\}$     & $\emptyset\mapsto \max(1,1)=1,\;\{3\}\mapsto \max(2,-\infty)=2$ \\
forget $3$ (root)   & $\emptyset$ & $\emptyset\mapsto \max(1,2)=2$ \\
\bottomrule
\end{tabular}
\end{center}

The root reports $2$. At ``intro $3$'' the state $\{3\}$
takes the child's $\emptyset\mapsto 1$ (which already encodes the choice ``$1$ in,
$2$ out'') and adds $w(3)=1$, giving $2$ -- the program has discovered $\{1,3\}$
without ever storing it explicitly. The $-\infty$ entries are the edges $12$ and
$23$ doing their job, forbidding adjacent vertices from being chosen together.

\subsection{Dominating set: a $3$-state DP}

A \emph{dominating set} $D\subseteq V$ is a set such that every vertex is in $D$
or adjacent to a vertex in $D$. Minimum dominating set is NP-hard. The bag state
now needs three colors per vertex, because a vertex's domination can be settled by
a neighbor that has not been processed yet.

For each vertex $v$ in the bag, assign one of:
\begin{itemize}
  \item \textcolor{sourceblue}{$\mathbf{1}$} (\emph{black}): $v\in D$.
  \item \textcolor{targetgreen}{$\mathbf{0}$} (\emph{white-dominated}): $v\notin
        D$ and $v$ is already dominated by some chosen vertex in $V_x$.
  \item \textcolor{quotientorange}{$\hat{\mathbf{0}}$} (\emph{white-pending}):
        $v\notin D$ and $v$ is \emph{not yet} dominated -- it is still waiting for
        a future neighbor (one introduced above $x$) to dominate it.
\end{itemize}
A coloring $\phi:B_x\to\{1,0,\hat 0\}$ is a state. Define $\mathrm{dp}[x][\phi]$ as
the minimum size of a set $D\subseteq V_x$ that (i) restricts to the black
vertices of $\phi$ on $B_x$, (ii) dominates every \emph{forgotten} vertex (every
vertex of $V_x\setminus B_x$), and (iii) dominates each white-$0$ bag vertex
already, while white-$\hat 0$ bag vertices are explicitly left undominated for
now. The transitions (minimizing):

\begin{itemize}
  \item \textbf{Leaf:} $\mathrm{dp}[x][\,]=0$ (empty coloring).
  \item \textbf{Introduce $v$:} $v$ has no forgotten neighbors yet. It either
        starts black ($+1$ to the size; this also upgrades any bag-neighbor that
        was $\hat 0$ to $0$) or stays out of $D$ at cost $0$. If it stays out,
        its color is forced by the bag: when some bag-neighbor is already black,
        $v$ is dominated at once and enters as white-$0$; when none is, $v$
        enters as white-pending $\hat 0$.
  \item \textbf{Forget $v$:} we drop $v$, but only from states where $v$ is
        \emph{not} pending: $v$ must be black or white-$0$ (it had its last chance
        to be dominated). We take the min over those child states. A child state
        with $v=\hat 0$ is discarded -- a pending vertex that is forgotten can
        never be dominated.
  \item \textbf{Join:} children $y,z$ with the same bag. Two colorings combine
        into $\phi$ when they agree on the black vertices, a bag vertex is
        white-$0$ in $\phi$ iff it is white-$0$ in \emph{at least one} child
        (dominated on either side counts), and white-$\hat0$ in $\phi$ iff it is
        $\hat0$ in \emph{both}. Add the sizes and subtract $\abs{\{v:\phi(v)=1\}}$
        for the double-counted black bag vertices.
\end{itemize}

The answer is $\mathrm{dp}[r][\,]$ at the empty root. Figure~\ref{ch24:fig-dom-states} illustrates how the three states interact across introduce, forget, and join steps.

A tiny trace makes the rules concrete. Take the path $a-b-c$ with the nice
decomposition: introduce $a$, introduce $b$, forget $a$, introduce $c$, forget
$c$, forget $b$. After introducing $b$ into $\{a,b\}$, the finite entries are
$(\hat0,\hat0)\mapsto 0$, $(1,0)\mapsto 1$, $(0,1)\mapsto 1$, and $(1,1)\mapsto 2$.
Forgetting $a$ discards the state where $a$ is $\hat 0$, leaving $b{=}1\mapsto 1$
and $b{=}0\mapsto 1$. Introducing $c$ and then forgetting it keeps only states
where $c$ ended up dominated: $b{=}1\mapsto 1$ (with $c$ white-$0$) and
$b{=}0\mapsto 2$ (with $c$ black). Forgetting $b$ takes the minimum, $1$, which is
the dominating set $\{b\}$.

\begin{figure}[H]
\centering
\begin{tikzpicture}[>=Stealth, every node/.style={font=\scriptsize},
  st1/.style={circle,draw=sourceblue,fill=sourceblue!18,thick,minimum size=6mm,inner sep=0pt,font=\scriptsize\bfseries},
  st0/.style={circle,draw=targetgreen,fill=targetgreen!15,thick,minimum size=6mm,inner sep=0pt,font=\scriptsize\bfseries},
  st0hat/.style={circle,draw=quotientorange,fill=quotientorange!15,thick,dashed,minimum size=6mm,inner sep=0pt,font=\scriptsize\bfseries},
  panel/.style={rounded corners=4pt,draw=gray!45,fill=gray!3,inner sep=6pt}
]
  % Panel 1: Introduce v=1
  \begin{scope}[xshift=0cm]
    \node[font=\footnotesize\bfseries,text=sourceblue!80!black] at (1.6,2.3) {1. Introduce $v$ as $\mathbf{1}$};
    \node[st0hat] (u1) at (0.5,1.3) {$\hat{0}$};
    \node[font=\tiny,left=1pt of u1] {$u$};
    \node[font=\tiny,text=gray!70!black] at (1.6,1.3) {child $y$};
    
    \draw[->,thick,structurepurple] (1.6,0.95) -- (1.6,0.45) node[midway,right=2pt,font=\tiny]{add $v{=}1$};

    \node[st0] (u2) at (0.5,0) {$0$};
    \node[font=\tiny,left=1pt of u2] {$u$};
    \node[st1] (v2) at (2.2,0) {$1$};
    \node[font=\tiny,right=1pt of v2] {$v$};
    \draw[thick,sourceblue] (u2) -- node[midway,above=1pt,font=\tiny]{dominates} (v2);
    \node[font=\tiny,align=center,text width=32mm] at (1.4,-0.85) {Neighbor $u$ upgrades\\from pending $\hat{0}$ to safe $0$};
  \end{scope}

  % Panel 2: Forget v
  \begin{scope}[xshift=4.6cm]
    \node[font=\footnotesize\bfseries,text=trapred!85!black] at (1.7,2.3) {2. Forget $v$ (Gatekeeper)};
    \node[st1] (f1) at (0.5,1.3) {$1$};
    \node[st0] (f0) at (1.7,1.3) {$0$};
    \node[st0hat] (fh) at (2.9,1.3) {$\hat{0}$};
    \draw[->,thick,targetgreen!80!black] (f1) -- (0.9,0.2);
    \draw[->,thick,targetgreen!80!black] (f0) -- (1.5,0.2);
    \draw[->,thick,trapred,dashed] (fh) -- (2.9,0.2) node[midway,right=2pt,font=\tiny\bfseries]{prune!};
    \node[rounded corners=2pt,draw=targetgreen,fill=targetgreen!10,inner sep=3pt] at (1.2,-0.1) {$\min(\text{cost})$};
    \node[text=trapred,font=\footnotesize\bfseries] at (2.9,-0.1) {$\times$};
    \node[font=\tiny,align=center,text width=34mm] at (1.7,-0.85) {Pending $\hat{0}$ cannot leave\\the bag undominated};
  \end{scope}

  % Panel 3: Join on vertex u
  \begin{scope}[xshift=9.4cm]
    \node[font=\footnotesize\bfseries,text=structurepurple!85!black] at (1.8,2.3) {3. Join Rule at $u \in B_x$};
    \node[st0] (jy) at (0.8,1.3) {$0$};
    \node[font=\tiny,above=1pt of jy] {in $y$};
    \node[st0hat] (jz) at (2.8,1.3) {$\hat{0}$};
    \node[font=\tiny,above=1pt of jz] {in $z$};
    \draw[->,thick,structurepurple] (jy) -- (1.6,0.25);
    \draw[->,thick,structurepurple] (jz) -- (2.0,0.25);
    \node[st0] (jx) at (1.8,-0.05) {$0$};
    \node[font=\tiny,align=center,text width=36mm] at (1.8,-0.85) {Logical OR: dominated in\\either subtree $\implies$ safe $0$};
  \end{scope}
\end{tikzpicture}
\caption{How the three dominating-set states -- $\mathbf{1}$ (in $D$), $\mathbf{0}$ (already dominated), and $\hat{\mathbf{0}}$ (still awaiting a dominator) -- flow up a nice tree decomposition. Introducing a black vertex $\mathbf{1}$ immediately upgrades its pending bag neighbors from $\hat{\mathbf{0}}$ to $\mathbf{0}$; forgetting a vertex discards any state where that vertex is still $\hat{\mathbf{0}}$; joining two subtrees takes the logical OR of domination on white vertices.}
\label{ch24:fig-dom-states}
\end{figure}

\begin{trapbox}
The classic bug is collapsing white-$0$ and white-$\hat0$ into one ``white''
state. Then the forget rule cannot tell whether a vertex about to vanish was ever
dominated, and the DP will happily produce sets that miss a vertex. The distinction between the two
kinds of white is the main subtlety of dominating set: ``not in $D$ and
already safe'' versus ``not in $D$ and still owing a dominator.'' A second trap is
the join rule -- domination is a logical OR across the two subtrees (dominated on
either side suffices), which is why white-$0$ in the parent needs white-$0$ in
\emph{just one} child, while pending requires pending in \emph{both}.
\end{trapbox}

\begin{complexitybox}
Per node, MWIS examines $2^{\abs{B_x}}\le 2^{k+1}$ states; introduce/forget cost
$\Oh(2^{k})$ and join costs $\Oh(2^{k})$ table combines. Over $\Oh(kn)$ nodes the
total is $\Oh(2^k\cdot kn)$, i.e.\ \emph{linear in $n$} for fixed $k$. Dominating
set has $3^{\abs{B_x}}$ states; a naive join is $\Oh(3^k\cdot 3^k)$ but with the
subset-convolution / fast-coloring trick it drops to $\Oh(3^k\cdot k)$ per join,
giving $\Oh(3^k\cdot kn)$ overall. In general a problem whose ``interface'' has
$c$ states per vertex runs in $\Oh(c^k\cdot\poly(k)\cdot n)$ time on a width-$k$
decomposition: single-exponential in treewidth, linear in the graph.
\end{complexitybox}

\section{Computing and Approximating Treewidth}
\label{ch24:compute}

So far we assumed a decomposition was handed to us. Where does it come from?
Deciding $\tw(\Gph)\le k$ is NP-complete (Arnborg--Corneil--Proskurowski, 1987),
but the parameter $k$ rescues us in three different ways.

\subsection{Exact, when treewidth is small: Bodlaender's theorem}

\begin{theorem}[Bodlaender, 1996]
\label{ch24:thm-bodlaender}
There is an algorithm that, given $\Gph$ and $k$, in time $2^{\Oh(k^3)}\cdot n$
either outputs a tree decomposition of width $\le k$ or reports $\tw(\Gph)>k$.
\end{theorem}

This is \emph{fixed-parameter tractable}: linear in $n$ for each fixed $k$, with
all the dependence shoved into the $2^{\Oh(k^3)}$ factor. The constant is far too large
in practice, so the result is mainly of theoretical interest: it shows that
``find a good decomposition'' is not the bottleneck for the FPT program. In
practice one uses heuristics or the approximation below.

\subsection{A clean constant-factor approximation}

A standard approximation routine, due in this form to Robertson--Seymour and sharpened by many,
builds a decomposition of width $\Oh(k)$ whenever $\tw\le k$, using
\emph{separators}. It relies on the fact that bounded-treewidth graphs
have balanced separators.

\begin{lemma}[Balanced separator from treewidth]
\label{ch24:lem-balanced-sep}
If $\tw(\Gph)\le k$ then for any weight on the vertices, there is a bag $B_x$ of
some optimal decomposition that is a $\tfrac12$-balanced separator: each component
of $\Gph\setminus B_x$ carries at most half the total weight. In particular every
$W\subseteq V$ admits a separator of size $\le k+1$ splitting $W$ in balanced
fashion.
\end{lemma}

\begin{proof}[Proof idea]
Root an optimal decomposition and push toward the heavier side: starting at the
root, repeatedly move to the child whose subtree carries more than half the
weight. This must stop at a node $x$ where no single child subtree exceeds half --
its bag $B_x$, of size $\le k+1$, is the balanced separator, because each
component of $\Gph\setminus B_x$ lies on one side of the tree split and inherits
the $\le\tfrac12$ bound from Lemma~\ref{ch24:lem-separator}.
\end{proof}

The recursive scheme then reads:

\begin{algorithm}[H]
\caption{\textsc{ApproxDecompose}$(U,W)$ -- decomposition of interior $U$ with boundary $W$}
\label{ch24:alg-approx}
\KwIn{an interior set $U$ and a boundary set $W$ with $U\cap W=\emptyset$ and $\abs W\le 3k+6$}
\KwOut{a tree decomposition of $\Gph[U\cup W]$ of width $\Oh(k)$ whose root bag contains $W$}
\If{$\abs{U\cup W}\le 4k+8$}{\Return single bag $U\cup W$\;}
find a separator $S$, $\abs S\le k+1$, that is $\tfrac{2}{3}$-balanced for $W$\;
\tcp{exists by Lemma~\ref{ch24:lem-balanced-sep} when $\tw\le k$, else report failure}
add one vertex of $U$ to $S$\tcp*{guarantees progress; now $\abs S\le k+2$}
\ForEach{component $\Gamma$ of $\Gph[U]\setminus S$}{
  $W_\Gamma\gets N(\Gamma)\cap (W\cup S)$\tcp*{boundary inherited by $\Gamma$}
  $\mathcal{T}_\Gamma\gets$\textsc{ApproxDecompose}$(\Gamma,\,W_\Gamma)$\;
}
create a root bag $W\cup S$ and attach each $\mathcal{T}_\Gamma$ as a child\;
\Return the assembled tree\;
\end{algorithm}

\begin{theorem}[$\Oh(1)$-approximation of treewidth]
\label{ch24:thm-approx}
If $\tw(\Gph)\le k$, Algorithm~\ref{ch24:alg-approx} (started with $U=V$ and
$W=\emptyset$) returns a tree decomposition of width $\le 4k+7=\Oh(k)$. With a
more careful separator routine, Bodlaender et al.\ (2016) obtain width $\le 5k+4$
in $2^{\Oh(k)}\cdot n$ time.
\end{theorem}

\begin{whybox}
Why does keeping $\abs W\le 3k+6$ work? Take a child component $\Gamma$. The set
$\Gamma$ together with its neighbors in $W\setminus S$ is connected and avoids
$S$, so it lies inside one component of $\Gph\setminus S$, which holds at most
$\tfrac23\abs W$ vertices of $W$. The child boundary is those neighbors plus at
most all of $S$:
\[
  \abs{W_\Gamma}\le \tfrac23(3k+6)+(k+2)=3k+6.
\]
So the bound is self-sustaining, and every bag $W\cup S$ has at most
$(3k+6)+(k+2)=4k+8$ vertices. The extra interior vertex placed in $S$ makes each
child's interior $\Gamma$ strictly smaller than $U$, so the recursion ends. The
factor lost (here, roughly $4$) comes from carrying the boundary plus the fresh
separator in each bag.
\end{whybox}

\subsection{Practical heuristics: elimination orderings}

Almost every practical solver uses a concrete equivalent of treewidth: \emph{vertex elimination}. Eliminate a vertex $v$ by deleting it and
adding edges to make all its former neighbors a clique (``fill-in''). An
\emph{elimination ordering} $\pi$ eliminates vertices one by one; its \emph{width}
is the largest neighbor-clique size encountered, minus one.

\begin{theorem}[Treewidth equals minimum elimination width]
\label{ch24:thm-elim}
$\tw(\Gph)$ equals the minimum over all orderings $\pi$ of the elimination width
of $\pi$. In addition, an ordering of width $k$ yields a tree decomposition of width
$k$ (and vice versa) in linear time.
\end{theorem}

\begin{proof}[Why the two notions coincide]
Given a width-$k$ decomposition, repeatedly take a leaf bag $B_x$ and eliminate a
vertex $v$ that appears \emph{only} in $B_x$ (it exists by \textbf{(C)} unless the
leaf duplicates its parent, in which case delete the leaf): $v$'s neighbors all
sit in $B_x$ by the separator property, so the fill-in clique has size
$\le\abs{B_x}-1\le k$. This produces an ordering of width $\le k$. Conversely,
given an elimination ordering, the bag of $v$ is $\{v\}\cup(\text{neighbors of }v
\text{ at its elimination, in the fill-in graph})$; arranging these bags by the
ordering and connecting each $v$ to the bag of its earliest-eliminated fill-in
neighbor yields a valid tree decomposition of the same width. The fill-in edges
are exactly the edges that \textbf{(E)} forces a bag to realize.
\end{proof}

Figure~\ref{ch24:fig-elimination} shows a single elimination step and why the order in which vertices are eliminated governs the fill-in.

\begin{figure}[H]
\centering
\begin{tikzpicture}[>=Stealth, every node/.style={font=\small},
  V/.style={circle,draw=nodeblue,fill=blue!5,thick,minimum size=6.5mm,inner sep=0pt},
  Velim/.style={circle,draw=trapred,fill=trapred!15,thick,minimum size=6.5mm,inner sep=0pt}
]
  % Panel 1: Eliminate d (degree 3, creates fill edge b-e)
  \begin{scope}[xshift=0cm]
    \node[font=\footnotesize\bfseries,text=trapred!85!black] at (1.6,2.8) {Eliminating $d$ (degree $3$)};
    \node[V] (b1) at (0.0,1.0) {$b$};
    \node[V] (c1) at (1.6,0.3) {$c$};
    \node[Velim] (d1) at (1.6,1.8) {$d$};
    \node[V] (e1) at (3.2,1.0) {$e$};
    \draw[thick] (b1)--(c1) (b1)--(d1) (c1)--(d1) (c1)--(e1) (d1)--(e1);
    \draw[very thick,trapred,dashed] (b1.south) .. controls (0.8,-0.55) and (2.4,-0.55) .. node[midway,below=2pt,font=\scriptsize]{fill-in $be$} (e1.south);
    \node[font=\scriptsize,align=center,text width=42mm] at (1.6,-1.25) {$N(d)=\{b,c,e\}$ lacked $be$;\\adds $1$ fill-in edge};
  \end{scope}

  % Divider arrow / vs
  \node[font=\footnotesize\bfseries,text=gray!70!black] at (4.8,0.9) {vs.};

  % Panel 2: Eliminate b first (degree 2, zero fill-in)
  \begin{scope}[xshift=6.2cm]
    \node[font=\footnotesize\bfseries,text=targetgreen!80!black] at (1.6,2.8) {Eliminating $b$ first (degree $2$)};
    \node[Velim] (b2) at (0.0,1.0) {$b$};
    \node[V] (c2) at (1.6,0.3) {$c$};
    \node[V] (d2) at (1.6,1.8) {$d$};
    \node[V] (e2) at (3.2,1.0) {$e$};
    \draw[thick] (b2)--(c2) (b2)--(d2) (c2)--(e2) (d2)--(e2);
    \draw[very thick,targetgreen!80!black] (c2)--(d2) node[midway,right=2pt,font=\tiny\bfseries]{clique};
    \node[font=\scriptsize,align=center,text width=44mm] at (1.6,-1.25) {$N(b)=\{c,d\}$ is already adjacent:\\$0$ fill-in, bag $\{b,c,d\}$ (width $2$)};
  \end{scope}
\end{tikzpicture}
\caption{Eliminating a vertex connects all of its current neighbors into a clique. On the core $\{b,c,d,e\}$ of our running graph, eliminating $d$ while $b,c,e$ remain requires adding the dashed fill-in edge $be$. Eliminating the degree-$2$ vertex $b$ first needs no fill-in because its neighbors $\{c,d\}$ already share an edge, immediately yielding the width-$2$ bag $\{b,c,d\}$ followed by $\{c,d,e\}$.}
\label{ch24:fig-elimination}
\end{figure}

Heuristics now pick the ordering greedily:

\begin{itemize}
  \item \textbf{Min-degree:} always eliminate a vertex of smallest current degree.
        Cheap, surprisingly effective.
  \item \textbf{Min-fill:} eliminate the vertex whose elimination adds the
        \emph{fewest} fill-in edges. Slower per step, usually tighter width.
\end{itemize}

Neither is guaranteed to be optimal -- finding a minimum-width elimination ordering is
equivalent to computing treewidth itself, and therefore NP-hard -- but on the structured instances that arise in practice (constraint
networks, probabilistic graphical models, register allocation) they routinely find
near-optimal decompositions.

\begin{labbox}
Run min-degree elimination on our seven-vertex running graph. Degrees: $\deg a=2,
\deg f=2, \deg g=2$ are the smallest. Eliminate $a$ first: its neighbors $\{b,c\}$
are already adjacent, so \emph{no} fill-in, clique size $2$. Then eliminate $f$
(neighbors $\{e,g\}$, already adjacent, no fill-in). Then $g$ (neighbor $\{e\}$,
clique size $1$). Now after removing $a,f,g$, vertex $d$ has neighbors $\{b,c,e\}$;
edges $bc$ and $ce$ exist, $be$ does not, so eliminating $d$ adds one fill edge
$be$, clique size $3$. The maximum clique size seen is $3$, width $3-1=2$. The
ordering certifies $\tw=2$, matching our hand decomposition.
\end{labbox}

\section{Treewidth and Fixed-Parameter Tractability}
\label{ch24:fpt}

The dynamic programs in this chapter are special cases of a general theorem,
which the next chapter develops in full.

\begin{theorem}[Courcelle, 1990]
\label{ch24:thm-courcelle}
Every graph property expressible in \emph{monadic second-order logic} (MSO$_2$:
quantification over vertices, edges, and sets thereof, with adjacency and
incidence) can be decided in time $f(\tw,\,\abs{\varphi})\cdot n$ on graphs given
with a width-$\tw$ tree decomposition, for some computable $f$.
\end{theorem}

Independent set (``there is a vertex set $I$ such that no edge has both endpoints
in $I$''), dominating set, $3$-colorability, Hamiltonicity -- all are
MSO$_2$-expressible, so all are linear-time on bounded-treewidth graphs by a
single meta-theorem. Our hand-built DPs are just Courcelle's automaton spelled out
for specific properties, usually with far better constants than the generic
construction.

\begin{connectionbox}
The bag-as-separator idea recurs throughout the book. The separator framing
connects forward to Chapter~27 (planar graphs have
treewidth $\Oh(\sqrt n)$, which is why ``divide along a small separator'' gives
subexponential planar algorithms). It connects to the dynamic programming on trees
you saw with heavy-light and centroid decompositions -- those are treewidth-$1$
decompositions specialized for fast path and subtree queries. And it connects to linear algebra: a tree
decomposition is exactly a sparse \emph{elimination ordering} for Gaussian
elimination, where the fill-in clique is the set of nonzeros created when you
pivot. Minimizing treewidth is minimizing fill-in, the central question of sparse
direct solvers and of nested dissection for solving $\Lap x=b$ on structured
graphs. When the companion \emph{Linear Algebra} volume discusses Cholesky fill-in and the
elimination tree, it is discussing treewidth in matrix language.
\end{connectionbox}

\begin{applicationbox}
Bounded treewidth matters in practice. \emph{Probabilistic inference}
in Bayesian networks and Markov random fields is exactly DP over a tree
decomposition of the model graph (the ``junction tree'' algorithm); its running
time is exponential in the treewidth, which is why practitioners work hard at
finding low-width elimination orders. \emph{Register allocation} models live
ranges as an interference graph; chordal interference graphs (treewidth $=$ clique
number $-1$) color optimally in polynomial time. \emph{Constraint satisfaction}
and SQL \emph{query optimization} use the treewidth (or its cousin, hypertree
width) of the constraint/join graph to bound evaluation cost. And many
\emph{road-network} and \emph{VLSI} instances have small treewidth in practice,
which is why preprocessing-based shortest-path engines use tree- and
path-decompositions.
\end{applicationbox}

\section{Pathwidth, Branchwidth, and the Wider Family}
\label{ch24:family}

Two close relatives deserve a mention. A \emph{path decomposition} is a tree
decomposition whose tree is a path; its minimum width is the \emph{pathwidth}
$\mathrm{pw}(\Gph)\ge\tw(\Gph)$. Pathwidth corresponds to a linear sweep with a
sliding window -- the spine decomposition of our running graph was already a path
decomposition, so it has pathwidth $2$ as well. \emph{Branchwidth} routes the
graph's edges through a ternary tree and differs from treewidth by at most a
constant factor: $\mathrm{bw}\le\tw+1\le\tfrac{3}{2}\,\mathrm{bw}$. These all
travel together, and the deepest fact in the area is the Grid Minor Theorem
(Robertson--Seymour): a graph of large treewidth contains a large grid as a minor,
and grids have large treewidth, so treewidth is, up to polynomial factors, the
same as ``largest grid you can find inside.'' This correspondence between a structural
obstruction (the grid) and a width parameter is the basis of
bidimensionality theory and subexponential FPT algorithms on minor-closed
classes.

\begin{connectionbox}
\textbf{What This Made Possible:}
Treewidth connects polynomial-time tree algorithms to NP-hard
problems on general graphs. When a graph has bounded treewidth ($k = \Oh(1)$), NP-hard
problems like Maximum Independent Set, 3-Coloring, Dominating Set, and Hamiltonian Cycle
become solvable in linear time $2^{\Oh(k)}\cdot n$. This gives exact solutions for
real-world graphs with tree-like structure, including control-flow graphs of structured
programs, telecommunication backbones, and RNA secondary structures.

\textbf{Limitations:}
\begin{itemize}
  \item \emph{Exponential explosion in width:} The running time scales as $2^{\Oh(k)}$
        or $3^{\Oh(k)}$. When $k \ge 30$, memory requirements and transition tables
        exceed practical computational limits.
  \item \emph{Meshes, grids, and dense graphs fail:} An $n\times n$ planar grid has
        treewidth $\Th(n)$, and general planar graphs can have treewidth $\Th(\sqrt{n})$.
        Expander graphs and dense cliques have treewidth $\Th(n)$, reducing the dynamic
        program to $2^{\Th(n)}$ brute force.
  \item \emph{Computing treewidth is NP-hard:} Determining whether $\tw(\Gph)\le k$ is
        NP-complete in general. While Bodlaender's algorithm can decide it in linear time
        for constant $k$, its constant factor $2^{\Om(k^3)}$ is astronomical for $k \ge 5$.
\end{itemize}

\textbf{Where We Go Next:}
Treewidth is only one structural parameter. When treewidth is large, can we parameterize
by solution size $k$, or preprocess the graph to remove easy parts?
In \textbf{Chapter~25 (Fixed-Parameter Tractability on Graphs)}, we develop the broader
theory of FPT: kernelization, crown decompositions, and bounded search trees.
For planar graphs where treewidth is $\Th(\sqrt{n})$, we develop topological tools in
\textbf{Chapter~26 (Planar Graphs)} and \textbf{Chapter~27 (Separators)}.
\end{connectionbox}

\section{Exercises}
\label{ch24:exercises}

\begin{exercise}[Treewidth of small graphs]
Compute the treewidth of: (a) the cycle $C_n$, (b) the complete bipartite graph
$K_{2,n}$, (c) the $3$-dimensional hypercube $Q_3$, (d) a single triangle with a
pendant path of length $5$ attached. Prove both the upper bound (exhibit a
decomposition) and the lower bound (exhibit a clique or use
Proposition~\ref{ch24:prop-trees}).
\end{exercise}

\begin{exercise}[The separator lemma, sharpened]
Prove a converse flavor of Lemma~\ref{ch24:lem-separator}: if every bag of a
purported decomposition acts as the stated separator and conditions
\textbf{(V)},\textbf{(E)} hold, must condition \textbf{(C)} hold? Either prove it
or give a counterexample.
\end{exercise}

\begin{exercise}[Implement MWIS]
Implement Algorithm~\ref{ch24:alg-mwis} end to end: a min-degree heuristic to get
an elimination ordering, conversion to a nice tree decomposition, and the DP. Test
it against brute force on random graphs of $12$ vertices, and on a $4\times4$ grid
(treewidth $4$). Report the empirical $2^k$ scaling.
\end{exercise}

\begin{exercise}[Reconstruct the witness]
Algorithm~\ref{ch24:alg-mwis} returns only the optimum \emph{value}. Augment it to
return an actual maximum-weight independent set by storing back-pointers and
walking down from the root. Prove your reconstruction is correct.
\end{exercise}

\begin{exercise}[Dominating set transitions in full]
Write out the complete introduce/forget/join tables for the dominating-set DP of
Section~\ref{ch24:dp}, including the bookkeeping that upgrades a bag-neighbor from
white-$\hat0$ to white-$0$ when a black vertex is introduced. Trace your rules on
the path $P_4$ and confirm the minimum dominating set has size $2$.
\end{exercise}

\begin{exercise}[Three-coloring]
Design a DP over a nice tree decomposition that decides $3$-colorability in
$\Oh(3^k\cdot\poly(k)\cdot n)$ time. The bag state is a coloring of $B_x$ with
$\{1,2,3\}$; state the introduce, forget, and join rules and prove the join rule
correct using the separator property.
\end{exercise}

\begin{exercise}[Break a fake decomposition]
Give three bags $B_1,B_2,B_3$ on a path $B_1-B_2-B_3$ that satisfy \textbf{(V)} and
\textbf{(E)} for some graph but violate \textbf{(C)}, and exhibit a concrete edge
of the graph for which the bag-as-separator lemma fails. This shows \textbf{(C)} is
not redundant.
\end{exercise}

\begin{exercise}[Elimination $=$ fill-in]
Prove the two directions of Theorem~\ref{ch24:thm-elim} in full detail for a
worked $6$-vertex graph: pick an ordering, compute the fill-in graph and its
width, then read off the corresponding tree decomposition and verify all three
conditions hold with the same width.
\end{exercise}

\begin{exercise}[Series-parallel graphs]
A graph is \emph{series-parallel} if it is built from a single edge by series and
parallel compositions. Prove every series-parallel graph has treewidth $\le 2$,
and conversely that every graph of treewidth $\le 2$ is a subgraph of a
series-parallel graph. (Hint: treewidth $\le 2$ graphs are exactly the subgraphs
of partial $2$-trees.)
\end{exercise}

\begin{exercise}[Pathwidth gap]
Exhibit a family of trees (treewidth $1$) whose pathwidth grows like
$\Theta(\log n)$. The complete binary tree of depth $d$ is the canonical example:
prove its pathwidth is roughly $\lceil d/2\rceil$. This shows pathwidth can be
unboundedly larger than treewidth.
\end{exercise}

\begin{exercise}[Grid lower bound]
Prove that the $n\times n$ grid has treewidth $\ge n$. (Hint: any balanced
separator of the grid has size $\ge n$, and Lemma~\ref{ch24:lem-balanced-sep}
turns small treewidth into small balanced separators; contrapose.)
\end{exercise}

\begin{exercise}[Treewidth under operations]
Show how treewidth changes under: adding one edge ($\tw$ rises by at most $1$),
contracting an edge ($\tw$ does not increase), and taking the disjoint union ($\tw$
is the max). Use these to bound the treewidth of $K_n$ minus a perfect matching.
\end{exercise}
